% generated from Paper 1/anccft.tex
\chapter[I. The Langlands--thermodynamic bridge over Bruhat--Tits trees]{Algorithmic non-commutative class field theory, I: the Langlands--thermodynamic bridge over Bruhat--Tits trees}\label{chap:I}
\chaptermark{I}
\begin{chapterabstract}
We open the non-abelian, algorithmic half of the programme: the localized Bost--Connes
formalism is transported from the one-dimensional valuation space to the Bruhat--Tits
tree of $\GL_2$ over a non-archimedean local field, and the arithmetic input is upgraded
from Frobenius classes in an abelian Galois group to Satake parameters of unramified
automorphic representations.

Section~\ref{I:sec:geometry} constructs the state space $\Omega$ (the space of lattice
chains, equivalently the tree boundary $\partial X\cong\PPone(F)$), the truncated Hecke
groupoids $\cG^{(N)}$ and their limit $\cG=\bigcup_N\cG^{(N)}$, and the tail equivalence
relation $\cR$ with its Hecke-degree cocycle; the groupoid is \'etale, amenable and
essentially principal, and its algebra is a Cuntz algebra locally and a Cuntz--Krieger
algebra $\cO_{A^{\mathrm{nb}}_Y}\simeq C(\partial X)\rtimes\Gamma$ globally.

Section~\ref{I:sec:thermo} inserts the Satake parameters into the Radon--Nikodym cocycle.
The partition function of the twisted system is \emph{exactly} the local $L$-factor,
$Z_p(\beta)=L_p(\beta,\pi_p)$, an identity of the Hecke recursion, and the abscissa of
convergence $\beta_c(\pi_p)=\log_q\max_i|\alpha_{p,i}|$ is the largest pole of $L_p$. Two
structural facts then follow, and they cut in opposite directions. Positivity
(Perron--Frobenius) shows that $\KMS$ \emph{states} see only the Perron branch of the
unramified spectrum, so cuspidal Satake data survive only as $\KMS$ \emph{functionals};
and the Kakutani dichotomy for the associated conformal measures has obstruction series
$\Xi_\beta(\pi,\pi';S)=\sum_{p\in S}p^{-\beta}\|\,\cdot\,\|^2$ governed by
Rankin--Selberg, giving a measure-theoretic form of strong multiplicity one. Combining
them yields the \emph{weight gap theorem}: equilibrium requires
$\beta>\beta_c(S)+w/2$ while the obstruction is confined to $\beta\le\beta_c(S)$, where
$w$ is the motivic weight; the two ranges meet if and only if $w=0$. The abelian theory
is the case $w=0$; Paper XXI's universal gap of width $1$ is the Eisenstein case $w=2$;
tempered weight two gives $w=1$ and a residual gap of exactly $\tfrac12$, the Ramanujan
exponent.

Section~\ref{I:sec:ktheory} is the algorithmic bridge, and it is where the higher-weight
arithmetic survives. Pimsner's exact sequence for the Hecke correspondence, the
non-commutative replacement of the equivariant Pimsner--Voiculescu sequence, identifies
$K_*$ with the kernel and cokernel of $1-u\,T_p+qu^2$ over $\Z[u^{\pm1}]$ --- the
\emph{Ihara module}, whose Fitting ideal is the Ihara zeta function by Bass's identity.
Specializing $u=1$ gives the Hecke--Laplacian $\Delta_p=(q+1)I-T_p$ and the structure
theorem
\[
K_0\cong\Z\oplus\Jac(Y),\qquad K_1\cong\Z,\qquad
|\Jac(Y)|=\frac1{|V|}\prod_f\bigl(p+1-a_p(f)\bigr)=\frac1{|V|}\prod_f\#E_f(\F_p),
\]
the torsion being computed exactly, as a group and not merely as an order, by the Smith
normal form of $\Delta_p$: Frobenius traces precipitate as $K$-theoretic torsion. By
contrast the boundary crossed product is arithmetically blind,
$K_0(C(\partial X)\rtimes\Gamma)\cong\Z^r\oplus\Z/(r-1)$ with $r$ the rank of $\Gamma$,
a theorem proved here by collapsing the reversal relation
$[f]+[\bar f]=c_{o(f)}$ to a single class of order $|E_Y|-|V_Y|$. Eleven independent
machine computations confirm both channels. Along Benjamini--Schramm convergent families
the torsion density converges exactly to Lyons' tree entropy,
$\frac1n\log|\Tors K_0|\to q\log q-\frac{q-1}2\log(q^2-1)$, and the Ramanujan bound
$2\log(1+\sqrt q)$, while unconditional, is strict and exceeds it by
$2q^{-1/2}+O(q^{-1})$ as $q\to\infty$ --- a Taylor expansion, not a fit --- because it
replaces a Kesten--McKay integral by its spectral edge. The Laplacian channel is realised as a
$\Z[u^{\pm1}]$-module and as the index of a Hecke--Dirac operator, not as a
$C^*$-algebra: the class $T_p-qI$ is a legitimate $KK$-class but not a correspondence,
and closing that gap is the paper's main open problem.
\end{chapterabstract}

\section*{Introduction, and the status of each statement}

The thirty-nine papers of \emph{Localized Bost--Connes systems} settle the
$\GL_1$ problem: the $\KMS_\beta$ simplex of $\cA_{K,S}$ is $\Prob(G_S/\Xi_\beta^\perp)$,
the obstruction group $\Xi_\beta$ is cut out by a Dirichlet series in the Frobenius
classes, the dichotomy underneath is Kakutani's, and the integral invariants are Smith
normal forms of $\Ns\pp\cdot 1-\Lambda^r\Pi$. Every ingredient of that theory has a
$\GL_2$ counterpart, and the counterparts are not formal analogies: valuations become
lattice chains, the profinite completion becomes the tree boundary, Frobenius classes
become Satake parameters, Dirichlet series become automorphic $L$-functions, and the
Smith normal form of a norm-minus-multiplication matrix becomes the Smith normal form of
a Hecke--Laplacian. This paper writes the dictionary and proves the entries that can be
proved.

It is written against a negative result of the earlier series, and we state at once how
it stands with respect to it. Paper XXI of that series proves that the untwisted $\GL_2$
system has a \emph{universal gap of width one}: the automorphic obstruction can be
non-trivial only for $\beta\le\beta_c(S)$, Gibbs states exist only for
$\beta>1+\beta_c(S)$, and no temperature does both. Nothing below contradicts that.
What changes here is \emph{where} the automorphic parameter is inserted --- into the
dynamics, not into the state --- and the effect of the change is quantitative and exact:
the gap is $w/2$, where $w$ is the motivic weight of the representation
(Theorem~\ref{I:thm:gap}). Paper XXI is the case $w=2$; the abelian theory of the monograph
is the case $w=0$, where the gap closes and the classification succeeds; tempered weight
two is $w=1$ and the residual gap is the Ramanujan exponent $\tfrac12$. The thermodynamic
channel therefore classifies exactly the weight-zero, that is Artin, part of the Langlands
correspondence. That is a sharp statement about what thermodynamics can and cannot do,
and it is the reason Section~\ref{I:sec:ktheory} exists: $K$-theory carries no positivity
constraint and no temperature, and it does see the higher-weight data.

\smallskip
\noindent\textbf{Status.} We label every numbered statement in one of three ways.
\emph{Theorem/Proposition}: proved here or an immediate assembly of quoted results, with
the assembly given. \emph{Verified}: an integral statement established by exact machine
computation on a stated finite list of cases, with a proof sketch but no proof; these are
marked $[\mathrm{V}]$ and the verification table is Table~\ref{I:tab:snf}; in the present
version no numbered statement carries this label, Theorem~\ref{I:thm:blind} having been
promoted to a proof, and Table~\ref{I:tab:snf} now serves as independent confirmation
rather than as evidence. \emph{Conjecture}:
neither. No statement below is asserted at a higher level than it is entitled to. In
particular Theorem~\ref{I:thm:kms} is a theorem, its extension to ramified and to $\GL_n$
data is Conjecture~\ref{I:conj:gln}, and the identification of the full extremal simplex
with the polar divisor of an $L$-function --- the form in which such a correspondence is
usually stated informally --- is \emph{false} as stated, for the reason isolated in
Theorem~\ref{I:thm:perron}.

\section{The geometric state space over Bruhat--Tits trees}\label{I:sec:geometry}

\subsection{The local field and its tree}

Throughout, $F$ is a non-archimedean local field with ring of integers $\cO$, maximal
ideal $\mathfrak{m}=(\varpi)$, uniformiser $\varpi$, normalised valuation
$v\colon F^\times\twoheadrightarrow\Z$, and residue field $k_F=\cO/\mathfrak{m}$ of
cardinality $q$. In the arithmetic applications $F=\Q_p$ and $q=p$. Write
$G=\GL_2(F)$, $K=\GL_2(\cO)$, $Z\cong F^\times$ for the centre.

\begin{definition}[Lattices and the tree]\label{I:def:tree}
An \emph{$\cO$-lattice} in $V=F^2$ is a free $\cO$-submodule $L\subset V$ of rank $2$.
Two lattices are \emph{homothetic}, $L\sim L'$, if $L'=cL$ for some $c\in F^\times$. The
Bruhat--Tits tree $X=X_F$ of $\PGL_2(F)$ has
\[
X^{(0)}=\{\text{lattices}\}/\!\sim,\qquad
[L]\text{ adjacent to }[L']\iff\exists\,L\in[L],\ L'\in[L']\ \text{with}\
\varpi L\subsetneq L'\subsetneq L .
\]
$X$ is a tree, it is $(q+1)$-regular, $G$ acts on it by $g\cdot[L]=[gL]$ with $Z K$ the
stabiliser of the standard vertex $x_0=[\cO^2]$, and the vertex set is
$X^{(0)}\cong G/ZK$.
\end{definition}

\begin{definition}[Boundary]\label{I:def:boundary}
A \emph{ray} from $x_0$ is a sequence $\xi=(x_0,x_1,x_2,\dots)$ of vertices with
$x_i\sim x_{i+1}$ and $x_{i+1}\neq x_{i-1}$ (\emph{non-backtracking}). The
\emph{boundary} $\partial X$ is the set of rays from $x_0$, topologised by the cylinder
sets $C(x_0,\dots,x_n)$. Then $\partial X$ is a compact totally disconnected space
without isolated points --- a Cantor set --- and the map sending a ray to the
corresponding line yields a $G$-equivariant homeomorphism
\[
\partial X\;\xrightarrow{\ \cong\ }\;\PPone(F)=\PPone(\cO)/\!\sim .
\]
The cylinder $C(x_0,\dots,x_n)$ has $q+1$ children at level $1$ and exactly $q$ children
at every later level.
\end{definition}

The two branching numbers $q+1$ and $q$ are the whole geometric content of the section,
and they are the reason the boundary and not the lattice space is the correct state
space; see Remark~\ref{I:rem:offbyone}.

\begin{definition}[Chain space]\label{I:def:chains}
Let
\[
\widetilde\Omega:=\Bigl\{(L_n)_{n\ge0}\ :\ L_0=\cO^2,\ L_{n+1}\subsetneq L_n,\
[L_n:L_{n+1}]=q\ \ \forall n\Bigr\},
\]
the space of \emph{descending lattice chains of step one}, with the inverse-limit
topology. Since the sublattices of index $q$ in a lattice $L$ are in bijection with
$\PPone(k_F)$, which has $q+1$ elements, $\widetilde\Omega$ is homeomorphic to the
boundary of the rooted $(q+1)$-ary tree, hence a Cantor set, and
$[L_0],[L_1],[L_2],\dots$ is a walk in $X$ issuing from $x_0$. The subspace
\[
\Omega:=\{(L_n)\in\widetilde\Omega:\ [L_{n+1}]\neq[L_{n-1}]\ \forall n\ge1\}
\]
of \emph{geodesic} chains is closed and $\Omega\cong\partial X\cong\PPone(F)$ by
Definition~\ref{I:def:boundary}. We call $\widetilde\Omega$ the \emph{Toeplitz} (or full)
state space and $\Omega$ the \emph{boundary} state space.
\end{definition}

\begin{remark}\label{I:rem:offbyone}
$\widetilde\Omega$ is the $\GL_2$ analogue of the space $Y_S$ of the abelian theory: its
level-$n$ set is the set of sublattices of index $q^n$ in $\cO^2$, of cardinality
$\sigma_1(q^n)=1+q+\dots+q^n$, and the associated Dirichlet series is the Solomon zeta
$\zeta_{M_2(\cO)}(\beta)=(1-q^{-\beta})^{-1}(1-q^{1-\beta})^{-1}$ of Paper XXI. The
passage $\widetilde\Omega\rightsquigarrow\Omega$ removes exactly the excess
$d_\pp=q+1>q$ that Paper XXI, Theorem 1.2 identifies as the source of its off-by-one:
on $\Omega$ the branching is $q$ per step and matches the norm.
\end{remark}

\subsection{The Hecke monoid and the truncated Hecke groupoids}

Let $\cH(G,K)=C_c(K\backslash G/K)$ be the spherical Hecke algebra with convolution, and
\[
T:=\bigl[K\begin{psmallmatrix}\varpi&0\\0&1\end{psmallmatrix}K\bigr],\qquad
R^{\pm1}:=\bigl[K\begin{psmallmatrix}\varpi^{\pm1}&0\\0&\varpi^{\pm1}\end{psmallmatrix}K\bigr],
\qquad
T_n:=\bigl[K\begin{psmallmatrix}\varpi^{n}&0\\0&1\end{psmallmatrix}K\bigr]\ (n\ge0),
\]
so that $\cH(G,K)=\C[T,R^{\pm1}]$ and the Satake isomorphism is
$\Sat\colon\cH(G,K)\xrightarrow{\cong}\C[\alpha_1^{\pm1},\alpha_2^{\pm1}]^{S_2}$ with
$\Sat(T)=q^{1/2}(\alpha_1+\alpha_2)$ and $\Sat(R)=\alpha_1\alpha_2$. Acting on functions
on $X^{(0)}$, $T$ is the adjacency (nearest-neighbour) operator and $R$ acts trivially.

\begin{lemma}[Hecke recursion]\label{I:lem:recursion}
On $X$ one has $T_1=T$, $T\,T_1=T_2+(q+1)T_0$ and, for $n\ge2$,
\[
T\,T_n=T_{n+1}+q\,T_{n-1},
\qquad\text{equivalently}\qquad
\sum_{n\ge0}T_n\,u^n=\bigl(1-Tu+qu^2\bigr)^{-1}
\]
as an identity of formal power series with coefficients in $\cH(G,K)$ normalised by
$\Sat(R)=q$. Consequently, for a character $\lambda$ of $\cH(G,K)$ with Satake
parameters $(\alpha_1,\alpha_2)$, $\alpha_1\alpha_2=q$,
\begin{equation}\label{I:eq:heckegen}
\sum_{n\ge0}\lambda(T_n)\,u^n=\frac1{(1-\alpha_1u)(1-\alpha_2u)} .
\end{equation}
\end{lemma}

The pencil $1-Tu+qu^2$ of Lemma~\ref{I:lem:recursion} is the single object that reappears
in every section below: in Section~\ref{I:sec:thermo} as the inverse of a partition
function, in Section~\ref{I:sec:ktheory} as the presentation matrix of a $K$-group.

\begin{definition}[Hecke monoid and degree]\label{I:def:monoid}
Let $S_\varpi:=\bigl(M_2(\cO)\cap G\bigr)/\cO^\times$, a monoid acting on lattices by
$L\mapsto gL$, and let $\deg\colon S_\varpi\to\Z_{\ge0}$, $\deg(g)=v(\det g)$. The
$K$-double cosets of $S_\varpi$ are indexed by elementary divisors, and
$\{\deg=n\}$ is the union of the double cosets of $T_a R^b$ with $a+2b=n$.
\end{definition}

\begin{definition}[Truncated Hecke groupoids]\label{I:def:groupoid}
Let $\Gamma\le\PGL_2(F)$ be a discrete, torsion-free, cocompact subgroup; such $\Gamma$
exist and are free of finite rank, and $Y:=\Gamma\backslash X$ is a finite
$(q+1)$-regular graph. Write $E_Y$ for its set of oriented edges, $\bar e$ for the
reversal of $e$, and
\[
A^{\mathrm{nb}}_{ef}=\begin{cases}1,&t(e)=o(f)\text{ and }f\neq\bar e,\\
0,&\text{otherwise},\end{cases}
\qquad e,f\in E_Y,
\]
for the \emph{non-backtracking} (edge) matrix. Let
$\Sigma_Y:=\{\omega\in E_Y^{\N}:A^{\mathrm{nb}}_{\omega_i\omega_{i+1}}=1\}$ with shift
$\sigma$. For $N\in\N$ set
\begin{equation}\label{I:eq:truncgroupoid}
\cG^{(N)}:=\bigl\{(\omega,i-j,\omega')\in\Sigma_Y\times\Z\times\Sigma_Y\ :\
0\le i,j\le N,\ \sigma^i\omega=\sigma^j\omega'\bigr\},
\end{equation}
and $\cG_\Gamma:=\bigcup_{N\ge0}\cG^{(N)}$, the \emph{Hecke groupoid}, with the
composition $(\omega,k,\omega')(\omega',k',\omega'')=(\omega,k+k',\omega'')$, the
topology generated by the sets
$Z(U,i,j,V)=\{(\omega,i-j,\omega'):\omega\in U,\ \omega'\in V,\ \sigma^i\omega=\sigma^j\omega'\}$
for $U,V$ cylinders, and the \emph{Hecke degree cocycle}
\[
c\colon\cG_\Gamma\longrightarrow\Z,\qquad c(\omega,k,\omega')=k .
\]
The same formula with $\Sigma_Y$ replaced by $\widetilde\Omega$ and $\sigma$ by the
one-sided shift on the full $(q+1)$-ary tree defines the \emph{Toeplitz Hecke groupoid}
$\widetilde\cG$.
\end{definition}

The truncation index $N$ is the number of Hecke moves allowed; $\cG^{(N)}$ is an open
subgroupoid, $\cG^{(N)}\subset\cG^{(N+1)}$, and $C^*(\cG_\Gamma)=\varinjlim_N C^*(\cG^{(N)})$.
This inductive limit is what carries the thermodynamic limit of
Section~\ref{I:sec:thermo}: a measure class on $\Omega$ is a compatible family of
quasi-invariant classes for the finite-level relations $\cR^{(N)}$ below.

\begin{definition}[Tail equivalence]\label{I:def:tail}
The \emph{tail (Hecke) equivalence relation} is the image of $\cG_\Gamma$ under
$(r,s)$,
\[
\cR:=\bigl\{(\omega,\omega')\in\Omega\times\Omega\ :\ \exists\, i,j\ge0,\
\sigma^i\omega=\sigma^j\omega'\bigr\},
\]
with the finite-level subrelations $\cR^{(N)}$ obtained by restricting $i,j\le N$, and
the \emph{degree-zero} subrelation $\cR_0:=\{(\omega,\omega'):\exists\,i,\ \sigma^i\omega=\sigma^i\omega'\}$.
A probability measure $\mu$ on $\Omega$ is \emph{$\beta$-quasi-invariant} for $\cR$ if
the measure class is preserved and the Radon--Nikodym cocycle of $\cG_\Gamma$ is
$e^{-\beta c}$; $\cR_0$ is a hyperfinite (AF) subrelation, and $\cR/\cR_0\cong\Z$ through
$c$.
\end{definition}

\begin{proposition}\label{I:prop:groupoid}
$\cG_\Gamma$ is a second countable, locally compact, Hausdorff, \'etale groupoid with
unit space $\Omega$; it is amenable, minimal and essentially principal, and
\begin{enumerate}
\item[(i)] $\cG_\Gamma$ is canonically isomorphic to the transformation groupoid
$\partial X\rtimes\Gamma$ under $\Omega\cong\partial X$, the Hecke degree $c$
corresponding to the Busemann (horocycle) cocycle of $\Gamma$ acting on $\partial X$;
\item[(ii)] $C^*(\cG_\Gamma)=C^*_r(\cG_\Gamma)$ is the Cuntz--Krieger algebra
$\cO_{A^{\mathrm{nb}}_Y}$, simple, purely infinite, nuclear, and Morita equivalent to
$C(\partial X)\rtimes\Gamma$;
\item[(iii)] in the rooted local model, $C^*(\widetilde\cG)\cong\cO_{q+1}$ and its
boundary (geodesic) reduction is stably isomorphic to $\cO_q$, so that
\[
K_0\bigl(C^*(\widetilde\cG)|_\Omega\bigr)\cong\Z/(q-1),\qquad
K_1\bigl(C^*(\widetilde\cG)|_\Omega\bigr)=0 .
\]
\end{enumerate}
\end{proposition}

\begin{remark}\label{I:rem:unitclass}
Part (iii) is the exact $\GL_2$ analogue of the unit-class torsion $\Z/(\Ns\pp^h-1)$ of
Paper XXXVIII: locally the boundary quotient contributes $\Z/(q-1)$ and nothing else. All
arithmetic beyond the residue cardinality is therefore invisible at a single place, and
must be sought either in the global assembly (Section~\ref{I:sec:thermo}) or in the
quotient graph $Y$ (Section~\ref{I:sec:ktheory}). This is the structural reason the
$\GL_2$ theory cannot be purely local, and it has no counterpart in the $\GL_1$ theory,
where the local boundary already carries the class group.
\end{remark}

\section{The thermodynamic--automorphic correspondence}\label{I:sec:thermo}

\subsection{Satake weights in the Radon--Nikodym cocycle}

Fix an unramified irreducible admissible representation $\pi_p$ of $G$ with Satake
parameters $(\alpha_{p,1},\alpha_{p,2})$ normalised arithmetically,
$\alpha_{p,1}\alpha_{p,2}=q^{w}$ with $w$ the motivic weight ($w=0$ for Artin
parameters, $w=k-1$ for weight-$k$ forms; $w=1$ in the classical weight-two case), and
local $L$-factor
\[
L_p(s,\pi_p)=\prod_{i=1,2}\bigl(1-\alpha_{p,i}\,q^{-s}\bigr)^{-1}.
\]

\begin{definition}[Satake potential and twisted dynamics]\label{I:def:dynamics}
Let $h_\pi\colon X^{(0)}\to\C^\times$ be a Hecke eigenfunction, $T h_\pi=\lambda_\pi h_\pi$
with $\lambda_\pi=\alpha_{p,1}+\alpha_{p,2}$, normalised by $h_\pi(x_0)=1$. Define the
\emph{Satake potential} on oriented edges and the associated \emph{Ruelle--Hecke transfer
operator} on $C(\Omega)$ by
\[
\varphi_\pi(e):=\log\frac{h_\pi(t(e))}{h_\pi(o(e))},\qquad
\bigl(\cL_{\beta,\pi}f\bigr)(\omega):=\sum_{\sigma(\eta)=\omega}
q^{-\beta\,c(\eta)}\,e^{\varphi_\pi(\eta_1)}f(\eta).
\]
The \emph{twisted dynamics} $\sigma^{(\pi)}$ on $C^*(\cG_\Gamma)$ is the one-parameter
group determined on $C_c(\cG_\Gamma)$ by
\[
\sigma^{(\pi)}_t(f)(\gamma)=e^{\,it\,(\beta\text{-independent})\,c_\pi(\gamma)}f(\gamma),
\qquad
c_\pi(\gamma):=c(\gamma)\log q-\varphi_\pi(\gamma),
\]
so that the Radon--Nikodym cocycle of a $\KMS_\beta$ state is
\begin{equation}\label{I:eq:RN}
\frac{d\mu\circ\gamma}{d\mu}(\omega)=e^{-\beta c_\pi(\gamma)}
=q^{-\beta c(\gamma)}\Bigl(\frac{h_\pi(t\gamma)}{h_\pi(o\gamma)}\Bigr)^{\beta}.
\end{equation}
\end{definition}

Equation~\eqref{I:eq:RN} is the precise sense in which the Satake parameter enters the
Radon--Nikodym derivative: it is the Doob $h$-transform of the geometric cocycle by the
Hecke eigenfunction. The untwisted system of Paper XXI is $h_\pi\equiv1$.

\subsection{The partition function is the local \texorpdfstring{$L$}{L}-factor}

\begin{theorem}[Partition function]\label{I:thm:partition}
Let $Z_p(\beta;\pi):=\sum_{n\ge0}\lambda_\pi(T_n)\,q^{-n\beta}$ be the partition function
of the truncated tower $(\cG^{(N)})_N$ with the dynamics of
Definition~\ref{I:def:dynamics}. Then the series converges absolutely if and only if
$\beta>\beta_c(\pi_p)$, where
\[
\beta_c(\pi_p):=\log_q\max_i|\alpha_{p,i}|,
\]
and in that range
\[
\boxed{\;Z_p(\beta;\pi)=\prod_{i=1,2}\bigl(1-\alpha_{p,i}q^{-\beta}\bigr)^{-1}
=L_p(\beta,\pi_p).\;}
\]
Moreover $\beta_c(\pi_p)$ is exactly the largest real part of a pole of
$s\mapsto L_p(s,\pi_p)$, so the critical inverse temperature of the twisted system is the
abscissa of convergence of the local $L$-factor. In particular
$\beta_c=w/2$ when $\pi_p$ is tempered $($Ramanujan$)$, and $\beta_c=w$ when $\pi_p$ is the
Eisenstein point $(\alpha_1,\alpha_2)=(1,q^{w})$.
\end{theorem}

\begin{proof}
Substitute $u=q^{-\beta}$ in \eqref{I:eq:heckegen} of Lemma~\ref{I:lem:recursion}; the radius
of convergence of the right-hand side is $\min_i|\alpha_{p,i}|^{-1}$. The poles of
$L_p(s,\pi_p)$ are at $q^{-s}=\alpha_{p,i}^{-1}$, i.e. at
$\Re s=\log_q|\alpha_{p,i}|$.
\end{proof}

\begin{remark}
For the trivial (Eisenstein) parameter $w=2$, $(\alpha_1,\alpha_2)=(1,q)$, one recovers
$Z_p(\beta)=(1-q^{-\beta})^{-1}(1-q^{1-\beta})^{-1}$, the Solomon zeta of $M_2(\cO)$ and
the local factor of $\zeta(\beta)\zeta(\beta-1)$: Theorem~\ref{I:thm:partition} contains
the partition function of the Connes--Marcolli $\GL_2$-system, and of Paper XXI, as its
$h_\pi\equiv1$ specialisation.
\end{remark}

\subsection{Equilibrium states, and what they can see}

\begin{theorem}[$\KMS$ classification, local]\label{I:thm:kms}
Let $Y=\Gamma\backslash X$ be connected and non-bipartite, and let
$A^{\mathrm{nb}}_{\beta,\pi}$ be the non-backtracking matrix weighted by
$e^{-\beta c_\pi}$ as in \eqref{I:eq:RN}. Then:
\begin{enumerate}
\item[(i)] $\KMS_\beta\bigl(C^*(\cG_\Gamma),\sigma^{(\pi)}\bigr)\neq\emptyset$ if and only
if the spectral radius satisfies $\rho\bigl(A^{\mathrm{nb}}_{\beta,\pi}\bigr)=1$, i.e.
if and only if the Bowen--Ruelle pressure equation $P(-\beta c_\pi)=0$ holds;
\item[(ii)] in that case the $\KMS_\beta$ simplex is a single point, the state
associated with the Perron eigenvector, and the corresponding measure $\mu_{\beta,\pi}$
is the unique $e^{-\beta c_\pi}$-conformal measure for $\cR$;
\item[(iii)] the associated von Neumann factor is of type $\mathrm{III}_{q^{-\beta}}$ if
$c_\pi(\cG_\Gamma)\subseteq\Z\log q$, and of type $\mathrm{III}_1$ if the closed group
generated by $e^{-\beta c_\pi(\cG_\Gamma)}$ is $\R_{>0}$;
\item[(iv)] for the untwisted potential the equation in $(\mathrm{i})$ reads
$q^{1-\beta}=1$, so the transition is at $\beta_c=1$ on the boundary state space
$\Omega$, against $\beta_c=2$ on the Toeplitz state space $\widetilde\Omega$.
\end{enumerate}
\end{theorem}

Part (iv) is the promised removal of Paper XXI's off-by-one: it is bought by the passage
from $\widetilde\Omega$ to $\Omega$, i.e. by the boundary quotient, and by nothing else.
The next theorem says what the price is.

\begin{theorem}[Perron obstruction: states cannot see cuspidal parameters]\label{I:thm:perron}
Let $\sigma^{(\pi)}$ be as in Definition~\ref{I:def:dynamics}. A $\KMS_\beta$ state exists
only if the eigenfunction $h_\pi$ can be chosen strictly positive on $X^{(0)}$. On a
$(q+1)$-regular tree a positive $T$-eigenfunction with eigenvalue $\lambda$ exists if and
only if $|\lambda|\ge2\sqrt q$. Consequently:
\begin{enumerate}
\item[(i)] if $\pi_p$ is tempered $(|\lambda_\pi|<2\sqrt q$, i.e.
$|\alpha_{p,i}|=q^{w/2})$, the twisted system has \emph{no} $\KMS_\beta$ state at any
$\beta$; the Satake data survive only as the $\KMS_\beta$ \emph{functional}
$\phi_{\beta,\pi}$ obtained by analytic continuation in $\lambda$, whose partition
function is $L_p(\beta,\pi_p)$ by Theorem~\ref{I:thm:partition};
\item[(ii)] the extremal $\KMS_\beta$ states of the untwisted system are parameterised by
the Perron branch alone --- the Eisenstein/complementary-series locus
$\lambda\ge2\sqrt q$ of the unramified spectrum --- and their partition functions are the
corresponding $L$-factors $L_p(\beta,1\boxplus|\cdot|^{w})$;
\item[(iii)] therefore the assertion ``the extremal $\KMS_\beta$ simplex is
parameterised by the poles of the automorphic $L$-function'' is true on the Eisenstein
locus and false on the cuspidal locus; the correct general statement replaces states by
functionals, and the simplex by the polar divisor of a meromorphic family.
\end{enumerate}
\end{theorem}

\begin{remark}[Thermodynamic ghost states]\label{I:rem:ghost}
We propose the name \emph{thermodynamic ghost states} for the non-positive continuations
$\phi_{\beta,\pi}$ of Theorem~\ref{I:thm:perron}(i): they are $\sigma^{(\pi)}$-$\KMS$
functionals with the correct modular property and the correct partition function
$L_p(\beta,\pi_p)$, but they are not states, having no positive extension to the algebra
and hence no GNS representation and no thermal equilibrium interpretation. The
terminology is borrowed, with the usual caution, from the negative-norm states of gauge
quantisation: as there, the ghosts are indispensable to the bookkeeping and absent from
the physical spectrum. The physical content of Theorem~\ref{I:thm:perron} is then this. The
metric geometry of the tree fixes a threshold, $2\sqrt q$, which is exactly the tempered
bound; a Satake parameter satisfying the Ramanujan estimate lies \emph{below} it, and a
positive Hecke eigenfunction --- the Gibbs weight of the would-be equilibrium --- then
simply does not exist on $X$. Cuspidality and thermal equilibrium are mutually exclusive
on the Bruhat--Tits tree. The arithmetic of the cuspidal spectrum is therefore not
suppressed but displaced: it survives as analytic data off the state space, and, as
Section~\ref{I:sec:ktheory} shows, as integral data in $K$-theory, where no positivity is
required.
\end{remark}

\begin{remark}
Theorem~\ref{I:thm:perron} is the operator-algebraic shadow of a familiar fact: positivity
of a state is a Perron--Frobenius condition, and Perron--Frobenius selects the trivial
representation. It is the same phenomenon that makes the constant function the only
positive $\ell^2$-eigenfunction on a finite graph, and the same one that in Paper XXVII
forces a $\KMS$ state to be tracial on the coefficient algebra and thereby erases inner
twists. Any programme that hopes to read cuspidal data off an equilibrium \emph{state}
must first defeat this, and we do not know how to.
\end{remark}

\subsection{The Kakutani obstruction and Rankin--Selberg}

We now globalise. Let $S$ be a set of rational primes, $\beta_c(S)$ the convergence
exponent of $\sum_{p\in S}p^{-\beta}$, and for unramified $\pi=\otimes_p\pi_p$ let
\[
\Omega_S:=\prod_{p\in S}\Omega_p,\qquad
\mu^\pi_\beta:=\bigotimes_{p\in S}\mu_{\beta,\pi_p},
\]
the product of the local conformal measures of Theorem~\ref{I:thm:kms}(ii), taken on the
locus where they exist, that is on the Eisenstein and complementary-series range
$|\lambda_p|\ge2\sqrt q$, or of their $h$-transform representatives there.

\begin{theorem}[Kakutani dichotomy for Satake data]\label{I:thm:kakutani}
For any two such $\pi,\pi'$ the measures $\mu^\pi_\beta$ and $\mu^{\pi'}_\beta$ are
either equivalent or mutually singular, and they are equivalent if and only if
\begin{equation}\label{I:eq:obstruction}
\Xi_\beta(\pi,\pi';S):=\sum_{p\in S}p^{-\beta}\,
\bigl|\widehat\lambda_p(\pi)-\widehat\lambda_p(\pi')\bigr|^2<\infty,
\qquad
\widehat\lambda_p:=\frac{\lambda_p}{2\sqrt p}=\cos\theta_p ,
\end{equation}
the equivalence of the Hellinger criterion with \eqref{I:eq:obstruction} holding with
absolute implied constants depending only on the temperedness bound. The series
\eqref{I:eq:obstruction} is the exact $\GL_2$ analogue of the obstruction series
$\sum_\pp\Ns\pp^{-\beta}|1-\chi(\Frob_\pp)|^2$ of the abelian theory, with the Frobenius
character replaced by the normalised Satake trace.
\end{theorem}

\begin{corollary}[Measure-theoretic strong multiplicity one]\label{I:cor:multone}
Let $\pi,\pi'$ be cuspidal automorphic representations of $\GL_2(\A_\Q)$, unramified
outside a finite set, and $S$ the set of all but finitely many primes, so
$\beta_c(S)=1$. Expanding \eqref{I:eq:obstruction} into Rankin--Selberg logarithms,
\[
\Xi_\beta(\pi,\pi';S)=\log L_S(\beta,\pi\times\widetilde\pi)
+\log L_S(\beta,\pi'\times\widetilde{\pi'})
-2\,\Re\log L_S(\beta,\pi\times\widetilde{\pi'})+O(1),
\]
so that, by the pole of $L(s,\pi\times\widetilde\pi)$ at $s=1$ and the regularity and
non-vanishing there of $L(s,\pi\times\widetilde{\pi'})$ for $\pi\not\cong\pi'$
$($Jacquet--Shalika$)$, one has $\Xi_1(\pi,\pi';S)=\infty$ and hence
\[
\pi\not\cong\pi'\ \Longrightarrow\ \mu^\pi_1\perp\mu^{\pi'}_1 .
\]
Distinct automorphic representations occupy mutually singular thermodynamic phases.
\end{corollary}

\begin{remark}
The Rankin--Selberg $L$-function appears here legitimately, and its appearance is not in
conflict with Paper XXI, which shows that the \emph{one-representation} obstruction
$\Xi_\beta(\rho,S)$ is linear in $\Tr\rho(\Frob_\pp)$ and therefore involves only the
standard $L$-function, the $\Sym^2$-term being spurious. The quantity
\eqref{I:eq:obstruction} is a Hellinger distance \emph{between two} representations; it is
genuinely quadratic, and the second moment of Satake traces is by definition
Rankin--Selberg.
\end{remark}

\subsection{The weight gap}

\begin{theorem}[Weight gap]\label{I:thm:gap}
Let $\pi$ have motivic weight $w$ at almost every place, i.e.
$|\alpha_{p,i}|=p^{w/2}$ for tempered $\pi_p$. Then, for every set $S$ of primes:
\[
\text{Gibbs states exist}\iff\beta>\beta_c(S)+\tfrac{w}{2},
\qquad
\Xi_\beta(\pi,\pi';S)=\infty\ \text{possible}\iff\beta\le\beta_c(S).
\]
The two ranges are disjoint with a gap of width exactly $w/2$, and they meet if and only
if $w=0$. Consequently:
\begin{enumerate}
\item[(i)] $w=0$ (Artin parameters, weight-one forms, the abelian case): the gap
closes, and one recovers the classification of the localized Bost--Connes theory,
namely $\KMS_\beta\cong\Prob(G_S/\Xi_\beta^\perp)$;
\item[(ii)] $w=1$ (tempered weight two): the residual gap is $\tfrac12$, the Ramanujan
exponent;
\item[(iii)] $w=2$ (the Eisenstein/untwisted $\GL_2$ system): the gap is $1$, which is
Paper XXI's universal off-by-one.
\end{enumerate}
\end{theorem}

\begin{proof}[Proof sketch]
By Theorem~\ref{I:thm:partition}, the equilibrium range is the domain of absolute
convergence of the Euler product of the local factors $L_p(\beta,\pi_p)$,
$p\in S$, whose abscissa is $\beta_c(S)+\beta_c(\pi_p)$; in the tempered case this equals
$\beta_c(S)+w/2$. For the obstruction range, note that
$|\widehat\lambda_p|\le1$, so the series \eqref{I:eq:obstruction} is bounded above
by $4\sum_{p\in S}p^{-\beta}$ and is therefore finite for every
$\beta>\beta_c(S)$. Disjointness and the width follow
\end{proof}

Theorem~\ref{I:thm:gap} is the main negative--positive result of the section. It says that
the thermodynamic channel is exactly a \emph{weight-zero} instrument: it classifies the
Artin part of the Langlands correspondence and is blind, by an amount equal to half the
motivic weight, to everything above it. This is a sharpening, not a refutation, of Paper
XXI: the twist moves the equilibrium threshold from $\beta_c(S)+1$ down to
$\beta_c(S)+w/2$, and for weight-one forms it closes the gap altogether.

\begin{conjecture}[$\GL_n$ and ramified data]\label{I:conj:gln}
Let $F$ be local, $G_n=\GL_n(F)$, $X_n$ the Bruhat--Tits building of $\PGL_n(F)$, and
$\pi_p$ unramified with Satake parameters $(\alpha_{p,1},\dots,\alpha_{p,n})$. Then the
partition function of the Hecke groupoid of the chamber boundary of $X_n$, twisted as in
Definition~\ref{I:def:dynamics}, equals $L_p(\beta,\pi_p)=\prod_i(1-\alpha_{p,i}q^{-\beta})^{-1}$,
the critical inverse temperature is $\max_i\log_q|\alpha_{p,i}|$, the weight gap theorem
holds verbatim with $w$ the motivic weight, and for ramified $\pi_p$ the same statements
hold with $\Omega$ replaced by the boundary of the building with level structure and
$L_p$ by the ramified local factor. For $n=2$ and unramified data this is
Theorems~\ref{I:thm:partition}--\ref{I:thm:gap}.
\end{conjecture}

\section{Algorithmic \texorpdfstring{$K$}{K}-theory via Hecke operators}\label{I:sec:ktheory}

Throughout this section $\Gamma$ is torsion-free and cocompact, $Y=\Gamma\backslash X$ is
a finite connected $(q+1)$-regular graph with $n=|V_Y|$ vertices and $m=|E_Y|$ unoriented
edges, so that
\[
r:=\rank\Gamma=m-n+1,\qquad m=\tfrac{(q+1)n}{2},\qquad r-1=m-n=\tfrac{(q-1)n}{2}.
\]
$T_p\in M_n(\Z)$ denotes the Hecke (adjacency) operator of $Y$ and
$\Delta_p:=(q+1)I_n-T_p$ its \emph{Hecke--Laplacian}. By Ihara's theorem and the
Eichler--Shimura and Jacquet--Langlands correspondences, when $\Gamma$ is the unit group
of an Eichler order the spectrum of $T_p$ is
\[
\Spec T_p=\{q+1\}\ \cup\ \{a_p(f)\ :\ f\ \text{a weight-two newform of the relevant level}\},
\]
so the entries of $\Delta_p$ are Frobenius traces in disguise. This section turns that
disguise into an algorithm.

\subsection{The Hecke correspondence and the Pimsner sequence}

\begin{definition}[Hecke correspondence]\label{I:def:corr}
Let $A:=C(V_Y)\cong\C^n$ and let $\cE_p$ be the $A$--$A$ correspondence with underlying
space $C(E_Y)$, right action $(\xi\cdot a)(e)=\xi(e)a(t(e))$, left action
$(a\cdot\xi)(e)=a(o(e))\xi(e)$, and inner product
$\langle\xi,\eta\rangle(v)=\sum_{t(e)=v}\overline{\xi(e)}\eta(e)$. Then $\cE_p$ is
finitely generated projective, full, with injective left action, and its class in
$KK^0(A,A)=M_n(\Z)$ is
\[
[\cE_p]=T_p .
\]
Let $\widetilde A:=C_0(V_Y\times\Z)$ with the degree shift $S$, and let
$\widetilde\cE_p$ be the corresponding correspondence over $\widetilde A$ recording the
valuation of the determinant, so that
$[\widetilde\cE_p]=u\,T_p$ on $K_0(\widetilde A)=\Z[u^{\pm1}]^n$, $u:=[S]$.
\end{definition}

\begin{theorem}[Hecke--Pimsner sequence; equivariant PV]\label{I:thm:pimsner}
For any $A$--$A$ correspondence $\cE$ as above, the Cuntz--Pimsner algebra $\cO_\cE$ sits
in the six-term exact sequence
\[
\begin{array}{ccccc}
K_0(A)&\xrightarrow{\ 1-[\cE]\ }&K_0(A)&\longrightarrow&K_0(\cO_\cE)\\
\uparrow& & & &\downarrow\\
K_1(\cO_\cE)&\longleftarrow&K_1(A)&\xleftarrow{\ 1-[\cE]\ }&K_1(A)
\end{array}
\]
which for $A=\C^n$ collapses to the short exact sequence
\begin{equation}\label{I:eq:pimsnershort}
0\longrightarrow K_1(\cO_\cE)\longrightarrow\Z^n
\xrightarrow{\ (1-[\cE])^{t}\ }\Z^n\longrightarrow K_0(\cO_\cE)\longrightarrow0 .
\end{equation}
When $\cE$ is the correspondence of a single automorphism, \eqref{I:eq:pimsnershort} is the
Pimsner--Voiculescu sequence; when $A$ carries an action of a finite group $C$ and $\cE$
is $C$-equivariant, the same sequence holds with $\Z^n$ replaced by the
$R(C)$-module $K_0^C(A)$ and $1-[\cE]$ by its $R(C)$-linear extension.
\end{theorem}

\begin{definition}[Ihara module]\label{I:def:ihara}
The \emph{Ihara module} of $(\Gamma,p)$ is the $\Z[u^{\pm1}]$-module
\[
M_\Gamma(p):=\Z[u^{\pm1}]^{\,n}\big/\bigl(1-u\,T_p+q\,u^2\bigr)\Z[u^{\pm1}]^{\,n},
\]
the presentation matrix being the Hecke pencil of Lemma~\ref{I:lem:recursion}.
\end{definition}

\begin{proposition}[Bass's identity; the two channels]\label{I:prop:bass}
With $A^{\mathrm{nb}}_Y$ as in Definition~\ref{I:def:groupoid},
\begin{equation}\label{I:eq:bass}
\det\bigl(I_{2m}-u\,A^{\mathrm{nb}}_Y\bigr)
=(1-u^2)^{\,m-n}\,\det\bigl((1+qu^2)I_n-u\,T_p\bigr),
\end{equation}
so the $0$-th Fitting ideal of $M_\Gamma(p)$ is generated by the reciprocal
$Z_Y(u)^{-1}(1-u^2)^{-(m-n)}$ of the Ihara zeta function of $Y$. Identity
\eqref{I:eq:bass} has been verified symbolically for all graphs of
Table~\ref{I:tab:snf}. Specialising $u=1$:
\[
\bigl(1-uT_p+qu^2\bigr)\big|_{u=1}=(q+1)I_n-T_p=\Delta_p .
\]
\end{proposition}

Proposition~\ref{I:prop:bass} isolates the two channels available for computation. The
\emph{boundary channel} uses $A^{\mathrm{nb}}_Y$ and computes the $K$-theory of
$C(\partial X)\rtimes\Gamma$ directly. The \emph{Laplacian channel} uses $\Delta_p$, that
is the $u=1$ specialisation of the Ihara module. They are related by \eqref{I:eq:bass},
and --- this is the point of the section --- they carry completely different arithmetic.

\subsection{The structure theorem}

\begin{theorem}[Hecke--Laplacian structure theorem]\label{I:thm:structure}
Let $\Delta_p=(q+1)I_n-T_p$ and let
$\SNF(\Delta_p)=\mathrm{diag}(d_1\mid d_2\mid\dots\mid d_{n-1},0)$ be its Smith normal
form over $\Z$. Then:
\begin{enumerate}
\item[(i)] $\ker\Delta_p=\Z$ and
$\coker\Delta_p\cong\Z\oplus\bigoplus_{i=1}^{n-1}\Z/d_i\Z
=\Z\oplus\Jac(Y)$,
where $\Jac(Y)$ is the Jacobian $($critical, sandpile$)$ group of $Y$;
\item[(ii)] $\coker\Delta_p$ and $\ker\Delta_p$ are the $u=1$ specialisation
\[
M_\Gamma(p)\otimes_{\Z[u^{\pm1}]}\Z_{u=1}\cong\Z\oplus\Jac(Y),
\qquad
\mathrm{Tor}_1^{\Z[u^{\pm1}]}\bigl(M_\Gamma(p),\Z_{u=1}\bigr)\cong\Z,
\]
of the Ihara module of Definition~\ref{I:def:ihara}, and they are the integral index
groups of the \emph{Hecke--Dirac operator} of $Y$: choosing an orientation $E_Y^+$ and
writing $\partial\colon\Z^{E_Y^+}\to\Z^{V_Y}$, $\partial f=t(f)-o(f)$, for the simplicial
boundary map, the self-adjoint operator
\[
\mathsf{D}:=\begin{pmatrix}0&\partial\\ \partial^{\,t}&0\end{pmatrix}
\quad\text{on}\quad \ell^2(V_Y)\oplus\ell^2(E_Y^+),
\qquad
\mathsf{D}^2=\Delta_p\oplus\partial^{\,t}\partial ,
\]
has $\ker\mathsf{D}=\Z\oplus H_1(Y;\Z)=\Z\oplus\Z^{\,r}$, and the integral cokernel of
the even part $\mathsf{D}^2|_{\ell^2(V_Y)}=\Delta_p$ is $\Z\oplus\Jac(Y)$. \emph{No
existence claim for a $C^*$-algebra is made here}; see Remark~\ref{I:rem:noalgebra};
\item[(iii)] the order of the torsion is a product of Frobenius data: by Kirchhoff's
matrix-tree theorem and Ihara's theorem,
\begin{equation}\label{I:eq:torsionorder}
\bigl|\Tors K_0(\cO)\bigr|=|\Jac(Y)|=\kappa(Y)
=\frac1n\prod_{\lambda\in\Spec T_p\setminus\{q+1\}}\bigl(q+1-\lambda\bigr)
=\frac1n\prod_f\bigl(p+1-a_p(f)\bigr),
\end{equation}
$f$ running over the weight-two newforms occurring in $Y$; and if $E_f$ is the elliptic
curve attached to $f$ by Eichler--Shimura, then $p+1-a_p(f)=\#E_f(\F_p)$, so
\[
\bigl|\Tors K_0(\cO)\bigr|=\frac1n\prod_f\#E_f(\F_p).
\]
\end{enumerate}
The Smith normal form computes not merely the order but the group: the elementary
divisors $d_1\mid\dots\mid d_{n-1}$ are the invariant factors of $\Jac(Y)$.
\end{theorem}

\begin{proof}[Proof sketch]
(i) is the standard theory of the critical group: $\Delta_p$ is the Laplacian of the
$(q+1)$-regular graph $Y$, symmetric with kernel the constants, and its cokernel splits
as $\Z\oplus K(Y)$ with $K(Y)$ the critical group. (ii) is
Theorem~\ref{I:thm:pimsner} applied with index map $\Delta_p$. In (iii), Kirchhoff gives
$\kappa(Y)=\frac1n\prod_{i\ge2}\mu_i$ for the non-zero Laplacian eigenvalues
$\mu_i=(q+1)-\lambda_i$, $|K(Y)|=\kappa(Y)$, and Ihara--Eichler--Jacquet--Langlands
identifies $\Spec T_p\setminus\{q+1\}$ with the $a_p(f)$. The last identity is the
definition of $a_p$.
\end{proof}

\begin{remark}[Bulk and boundary: what is, and is not, an algebra]\label{I:rem:noalgebra}
Since $A=\C^n$ we have $KK^0(A,A)\cong M_n(\Z)$, so
\[
x_p:=[\cE_p]-q\cdot 1_A=T_p-qI_n\ \in\ KK^0(A,A)
\]
is a bona fide Kasparov class, and $1-x_p=\Delta_p$. What fails is not the class but its
\emph{realisation}: a class is the class of a $C^*$-correspondence only if its matrix has
non-negative entries, and $T_p-qI_n$ has $-q$ on the diagonal. Hence $x_p$ is a formal
difference of correspondences, no Cuntz--Pimsner algebra $\cO_\cE$ has index map
$\Delta_p$, and we assert none. The asymmetry between the two channels is therefore not
only arithmetic but ontological:
\begin{center}\small
\begin{tabular}{lll}
\toprule
& realisation & invariant\\
\midrule
boundary channel & a $C^*$-algebra, $\cO_{A^{\mathrm{nb}}_Y}\simeq C(\partial X)\rtimes\Gamma$
& $\Z^r\oplus\Z/(r-1)$, no arithmetic\\
bulk (Laplacian) channel & a $\Z[u^{\pm1}]$-module and an index & $\Z\oplus\Jac(Y)$, all the arithmetic\\
\bottomrule
\end{tabular}
\end{center}
Everything asserted in Theorem~\ref{I:thm:structure} is unconditional at the level of the
Ihara module, of the index of $\mathsf{D}$, and of the integers; what is missing is an
honest $C^*$-algebra whose $K_0$ is $\Z\oplus\Jac(Y)$. Producing one --- equivalently,
realising the virtual correspondence $x_p$ by a genuine, possibly non-unital or
$\Z/2$-graded, construction --- is the main open problem of this paper.

There is a purely algebraic first step, and it isolates exactly where the difficulty
lives. The Ihara pencil $1-uT_p+qu^2$ of Definition~\ref{I:def:ihara} is a quadratic
matrix polynomial, and it linearises in the standard way over $A^{\oplus2}$,
$A=\C^n$: setting
\[
C:=\begin{pmatrix}T_p&-qI_n\\ I_n&0\end{pmatrix}\ \in\ M_{2n}(\Z),
\qquad
\det\bigl(I_{2n}-uC\bigr)=\det\bigl(I_n-uT_p+qu^2I_n\bigr),
\]
a companion-matrix identity verified directly by block Gaussian elimination on
$I-uC$. Thus $C$ is a single, honest $\Z$-linear class in $KK^0(A^{\oplus2},A^{\oplus2})$
representing the whole Ihara pencil at once, with $[\cE_p]=T_p$ visible in its upper
left block; the obstruction of the previous paragraph has been moved, not removed, since
$C$ still has the non-positive block $-qI_n$ in its upper right corner and so is still
not the matrix of a $C^*$-correspondence over $A^{\oplus2}$.

The problem left for the sequel is accordingly sharp: \emph{dilate} $C$, in the sense of
dilation theory, to a genuinely non-negative matrix over a larger vertex algebra
$\widetilde A\supset A^{\oplus2}$ --- for instance by the Drinen--Tomforde
desingularisation of a directed graph carrying $C$ as a signed incidence matrix, with the
negative entries reinterpreted as sinks absorbing mass at rate $q$ per vertex --- in such
a way that the resulting Cuntz--Pimsner (or graph) algebra $\cO_{\widetilde C}$ receives
$A^{\oplus2}$ as a full corner and has $K_0(\cO_{\widetilde C})\cong\Z\oplus\Jac(Y)$ by
Morita reduction back to $\Delta_p$. We do not carry this out here; it is the intended
content of a sequel to this paper, and the companion linearisation above is offered as
its starting point.
\end{remark}

\begin{theorem}[Boundary channel is arithmetically blind]\label{I:thm:blind}
Let $Y$ be any finite connected $(q+1)$-regular graph and $\Gamma$ the corresponding
torsion-free cocompact lattice, $r=\rank\Gamma=m-n+1\ge2$. Then
\[
K_0\bigl(C(\partial X)\rtimes\Gamma\bigr)\cong\Z^{\,r}\oplus\Z/(r-1)\Z,
\qquad
K_1\bigl(C(\partial X)\rtimes\Gamma\bigr)\cong\Z^{\,r},
\qquad r=\rank\Gamma=\tfrac{(q-1)n}{2}+1 .
\]
The torsion of the boundary quotient therefore depends only on the Euler characteristic
$r-1=|E_Y|-|V_Y|$ of $Y$, hence only on the isomorphism class of $\Gamma$ as an abstract
free group: \emph{no Hecke eigenvalue is visible in the boundary channel.}
\end{theorem}

\begin{proof}
Write $A^{\mathrm{nb}}=T^{t}S-J$, where $S,T\in M_{n\times2m}(\Z)$ are the incidence
matrices $S_{v,e}=[o(e)=v]$, $T_{v,e}=[t(e)=v]$ and $J$ is the edge-reversal involution,
and put $M:=I_{2m}-A^{\mathrm{nb}}=I+J-T^{t}S$. Since $\SNF(M)=\SNF(M^{t})$ it suffices
to compute $\coker M=\Z^{E_Y}/M\Z^{E_Y}$, and by Theorem~\ref{I:thm:pimsner} this is
$K_0$, with $K_1=\ker M^t$.

Reading off the columns of $M$, the presentation of $\coker M$ has generators $[e]$,
$e\in E_Y$, and one relation for each $f\in E_Y$,
\begin{equation}\label{I:eq:blindrel}
[f]+[\bar f]=c_{o(f)},\qquad\text{where}\quad c_v:=\sum_{t(e)=v}[e].
\end{equation}
Applying \eqref{I:eq:blindrel} to $f$ and to $\bar f$ gives $c_{o(f)}=c_{t(f)}$ for every
edge, so $c_v=:c$ is independent of $v$ because $Y$ is connected. Fix an orientation
$E_Y^+$, use $[\bar f]=c-[f]$ to eliminate the reversed generators, and let
$d^-(v):=\#\{f\in E_Y^+:o(f)=v\}$. Expanding $c_v$ in the surviving generators turns the
$n$ remaining relations into
\[
R_v:\qquad \bigl(\partial^{\,t}\text{-column at }v\bigr)\ :\
\sum_{f\in E_Y^+,\,t(f)=v}[f]-\sum_{f\in E_Y^+,\,o(f)=v}[f]
=\bigl(1-d^-(v)\bigr)\,c .
\]
Hence $\coker M\cong\coker B$ for the $(m+1)\times n$ integer matrix
\[
B=\begin{pmatrix}\partial^{\,t}\\ -\lambda^{t}\end{pmatrix},
\qquad \lambda_v:=1-d^-(v),
\]
on generators $E_Y^+\cup\{c\}$. Projecting away the last coordinate gives the exact
sequence
\[
0\longrightarrow\langle c\rangle\longrightarrow\coker B
\longrightarrow\coker\bigl(\partial^{\,t}\bigr)=H^1(Y;\Z)\cong\Z^{\,r}
\longrightarrow0 ,
\]
$H^1$ of a graph being free of rank $r$. Finally $t\,c=0$ in $\coker B$ if and only if
$(0,\dots,0,t)\in\operatorname{im}B$, i.e. if and only if there is $x\in\Z^{V_Y}$ with
$\partial^{\,t}x=0$ and $-\lambda\cdot x=t$; the first condition forces $x=a\mathbf1$,
and then
\[
-\lambda\cdot x=-a\sum_{v}\bigl(1-d^-(v)\bigr)=-a\,(n-m)=a\,(r-1),
\]
so $c$ has order exactly $r-1$. The quotient $\Z^r$ being free, the sequence splits and
$\coker M\cong\Z^{r}\oplus\Z/(r-1)$; $\ker M^t$ is free of rank $2m-\rank M=r$.
\end{proof}

\begin{remark}[Where the spectrum goes]\label{I:rem:wherespectrum}
The proof shows exactly how the Hecke data disappear, and it is worth recording that the
mechanism is \emph{not} the one suggested by Bass's identity \eqref{I:eq:bass}: at $u=1$
both factors of \eqref{I:eq:bass} vanish, so \eqref{I:eq:bass} degenerates to $0=0$ and
carries no integral information there. What kills the spectrum is the local relation
\eqref{I:eq:blindrel}: it forces the $n$ a priori distinct classes $c_v$ to collapse to a
single class $c$, after which the only surviving numerical invariant is
$\sum_v(1-d^-(v))=n-m$. The orientation-dependent quantities $d^-(v)$ do not survive
individually --- only their sum does --- and the sum is the Euler characteristic. The
adjacency matrix $T_p$ never enters the computation at all.
\end{remark}

\begin{remark}[The dichotomy]\label{I:rem:dichotomy}
Theorems~\ref{I:thm:structure} and \ref{I:thm:blind} together are the main structural finding
of this paper. The boundary quotient --- the natural non-commutative object, the simple
purely infinite Cuntz--Krieger algebra of the tail relation --- has $K$-theory of purely
topological content. The arithmetic sits one specialisation away, in the Hecke pencil
$1-uT_p+qu^2$ at $u=1$; equivalently, the arithmetic is the \emph{first-order} datum of
the Ihara zeta function at its trivial zero. This is the exact non-commutative analogue
of the classical statement that $\zeta_Y(u)$ has a trivial factor $(1-u^2)^{m-n}$ and
that the interesting arithmetic is in the reduced factor.
\end{remark}

\begin{corollary}[Frobenius traces as torsion]\label{I:cor:frobenius}
Let $Y$ carry a single non-trivial Hecke eigenvalue $a_p$ with multiplicity $n-1$. Then
$\Jac(Y)\cong(\Z/(p+1-a_p))^{\,n-2}$ whenever the Laplacian is of the form
$cI-J$, and in general $|\Jac(Y)|=\frac1n\prod_f\#E_f(\F_p)$ exactly. In particular the
Frobenius trace $a_p$ is recovered from $K_0$ by
\[
a_p=p+1-\bigl(n\cdot|\Tors K_0|\bigr)^{1/(n-1)}
\]
in the eigenvalue-homogeneous case, and in general from the elementary divisors by the
inverse of the Smith algorithm.
\end{corollary}

\begin{corollary}[Ramanujan bound; strict, and never attained]\label{I:cor:ramanujan}
If $Y$ is Ramanujan, i.e. $|\lambda|\le2\sqrt q$ for every
$\lambda\in\Spec T_p\setminus\{\pm(q+1)\}$, then
\[
\bigl|\Tors K_0\bigr|\ \le\ \frac1n\bigl(q+1+2\sqrt q\bigr)^{\,n-1}
=\frac1n\bigl(1+\sqrt q\bigr)^{2(n-1)},
\]
so that, with $B(q):=2\log(1+\sqrt q)$,
\[
\limsup_{n\to\infty}\frac1n\log\bigl|\Tors K_0\bigr|\ \le\ B(q).
\]
The bound is unconditional on the Ramanujan locus, but it is \emph{never} an equality,
and the deficit is exponential in $n$: see Theorem~\ref{I:thm:bstorsion} and
Remark~\ref{I:rem:bsnumerics}.
\end{corollary}

\begin{theorem}[Benjamini--Schramm asymptotic torsion]\label{I:thm:bstorsion}
Let $(Y_j)_{j\ge1}$ be finite connected $(q+1)$-regular graphs, $n_j=|V_{Y_j}|\to\infty$,
which Benjamini--Schramm converge to the $(q+1)$-regular tree $T_{q+1}$ --- for instance
any family of girth tending to infinity, such as the Lubotzky--Phillips--Sarnak graphs
$X^{p,\ell}$ at fixed $p=q$. Let $\cO_j$ be the associated Hecke--Laplacian invariant of
Theorem~\ref{I:thm:structure}. Then the torsion density converges, and its limit is Lyons'
tree entropy of $T_{q+1}$:
\begin{equation}\label{I:eq:bslimit}
\lim_{j\to\infty}\frac1{n_j}\log\bigl|\Tors K_0(\cO_j)\bigr|
\;=\;h(q)\;:=\;q\log q-\frac{q-1}{2}\log\bigl(q^2-1\bigr)
\;=\;\log\frac{q^{\,q}}{(q^2-1)^{(q-1)/2}} ,
\end{equation}
equivalently $h(q)=\int_{-2\sqrt q}^{2\sqrt q}\log\bigl(q+1-x\bigr)\,d\mu_{\mathrm{KM}}(x)$
against the Kesten--McKay measure
\[
d\mu_{\mathrm{KM}}(x)=\frac{(q+1)\sqrt{4q-x^2}}{2\pi\bigl((q+1)^2-x^2\bigr)}\,dx ,
\]
the spectral measure of the simple random walk on $T_{q+1}$ \cite{Kesten1959,McKay1981}.
Moreover $h(q)<B(q)$ for every $q$, with
\[
\frac{B(q)}{h(q)}=2.1061,\ 1.3889,\ 1.2340,\ 1.0957,\ 1.0089
\quad\text{for}\quad q=2,5,9,29,999.
\]
\end{theorem}

\begin{corollary}[The deficit, to two orders]\label{I:cor:deficit}
As $q\to\infty$,
\begin{equation}\label{I:eq:deficit}
B(q)-h(q)\;=\;2q^{-1/2}-\tfrac32q^{-1}+O\bigl(q^{-3/2}\bigr) ,
\end{equation}
an asymptotic equality, not a numerical fit. Consequently, on graphs with $n$
vertices, the Ramanujan bound overstates $|\Tors K_0|$ by the factor
\[
\exp\bigl((B(q)-h(q))\,n\bigr)
\;=\;
\exp\bigl(2n\,q^{-1/2}+O(n\,q^{-1})\bigr) ;
\]
at $q=2$, $n=2000$ this factor is $e^{1852}$.
\end{corollary}

\begin{proof}
Write $h(q)=\log q-\tfrac{q-1}2\log(1-q^{-2})$ and expand
$-\log(1-q^{-2})=q^{-2}+\tfrac12q^{-4}+O(q^{-6})$, so
\[
h(q)=\log q+\tfrac{q-1}2\bigl(q^{-2}+O(q^{-4})\bigr)
=\log q+\tfrac12q^{-1}-\tfrac12q^{-2}+O(q^{-3}) .
\]
Write $B(q)=\log q+2\log(1+q^{-1/2})$ and expand
$\log(1+q^{-1/2})=q^{-1/2}-\tfrac12q^{-1}+\tfrac13q^{-3/2}+O(q^{-2})$, so
\[
B(q)=\log q+2q^{-1/2}-q^{-1}+\tfrac23q^{-3/2}+O(q^{-2}) .
\]
Subtracting, the $\log q$ terms cancel and
$B(q)-h(q)=2q^{-1/2}-q^{-1}-\tfrac12q^{-1}+O(q^{-3/2})
=2q^{-1/2}-\tfrac32q^{-1}+O(q^{-3/2})$, which is \eqref{I:eq:deficit}. (A direct
computer-algebra expansion confirms the next two orders as well, and that the
$q^{-2}$ coefficient vanishes identically --- the $-\tfrac12q^{-2}$ term present in
$B(q)$ and in $h(q)$ separately cancels exactly on subtraction:
$B(q)-h(q)=2q^{-1/2}-\tfrac32q^{-1}+\tfrac23q^{-3/2}+\tfrac25q^{-5/2}+O(q^{-3})$.) The
numerical values of Remark~\ref{I:rem:bsnumerics} and of the ratios above were computed
independently of this expansion and agree with it: at $q=10^4$, \eqref{I:eq:deficit}
predicts $0.02-0.00015=0.01985$ against the quadrature value $0.019852$.
\end{proof}

\begin{proof}[Proof]
By Theorem~\ref{I:thm:structure}(iii), $|\Tors K_0(\cO_j)|=\kappa(Y_j)$, so
\eqref{I:eq:bslimit} is Lyons' theorem on the asymptotics of the number of spanning trees
along Benjamini--Schramm convergent sequences of bounded degree \cite{Lyons2005},
applied to the limit $T_{q+1}$, together with Kirchhoff's formula in the form
$\frac1n\log\kappa=\frac1n\bigl[\sum_{i\ge2}\log(q+1-\lambda_i)-\log n\bigr]$ and the
convergence of the empirical spectral measures of $T_p$ to $\mu_{\mathrm{KM}}$. The
closed form and the integral representation of $h(q)$ agree to six decimals for
$q+1\in\{3,4,6,8,14,30\}$ by numerical quadrature. Strictness of $h<B$ is the strict
concavity of $\log$ together with the fact that $\mu_{\mathrm{KM}}$ is not a point mass
at the spectral edge.
\end{proof}

\begin{remark}[Why the Ramanujan bound is the wrong instrument]\label{I:rem:bsnumerics}
Corollary~\ref{I:cor:ramanujan} replaces the whole spectrum by its extreme point
$2\sqrt q$. But $\mu_{\mathrm{KM}}$ vanishes like a square root at $\pm2\sqrt q$ and puts
almost no mass there, so the bound is a supremum where the truth is an integral; this,
and not any slack in the choice of family, is the source of the exponential deficit of
Theorem~\ref{I:thm:bstorsion}.

The deeper point is that the Ramanujan property is neither necessary nor sufficient for
the limit \eqref{I:eq:bslimit}, and that the invariant which does control the torsion
density is Benjamini--Schramm convergence. Both halves are visible in exact computation.
Below, $\lambda_2$ is the largest non-trivial eigenvalue in absolute value,
$\mathrm{KS}$ the Kolmogorov--Smirnov distance from the empirical spectral distribution
of $T_p$ to $\mu_{\mathrm{KM}}$, and the last two columns compare the measured torsion
density with $h(q)$ and with the Ramanujan bound $B(q)$.
\begin{center}\small
\renewcommand{\arraystretch}{1.15}\setlength{\tabcolsep}{5pt}
\begin{tabular}{lccccccc}
\toprule
$Y$ & $n$ & $q$ & $\lambda_2$ & Ram.\ ? & $\mathrm{KS}$
& ratio to $h(q)$ & ratio to $B(q)$\\
\midrule
random $3$-regular   & $200$  & $2$ & $2.8315$ & no  & $0.0142$ & $0.9774$ & $0.4641$\\
random $3$-regular   & $1000$ & $2$ & $2.8261$ & yes & $0.0032$ & $0.9937$ & $0.4719$\\
random $3$-regular   & $2000$ & $2$ & $2.8239$ & yes & $0.0021$ & $0.9964$ & $0.4731$\\
$\mathrm{PG}(2,5)$   & $62$   & $5$ & $2.2361$ & yes & $0.2661$ & $0.9669$ & $0.6962$\\
LPS $X^{5,13}$       & $2184$ & $5$ & $4.2497$ & yes & $0.0375$ & $0.9982$ & $0.7187$\\
LPS $X^{5,17}$       & $4896$ & $5$ & $4.3089$ & yes & $0.0307$ & $0.9991$ & $0.7194$\\
\bottomrule
\end{tabular}
\end{center}
Two rows falsify the spectral-gap reading. The random $3$-regular graph on $200$
vertices has $\lambda_2=2.8315>2\sqrt2=2.8284$ and is therefore \emph{not} Ramanujan,
yet its torsion density is already within $2.3\%$ of $h(2)$. The incidence graph of
$\mathrm{PG}(2,5)$ is Ramanujan with room to spare, $\lambda_2=\sqrt5=2.2361$ against
the permitted $2\sqrt5=4.4721$, yet it has girth $6$, its spectral distribution is far
from Kesten--McKay ($\mathrm{KS}=0.27$), and its torsion density falls $3.3\%$
\emph{below} $h(5)$. The $\mathrm{KS}$ column and the $h$-ratio column move together
throughout; the $\lambda_2$ column does not track either.

One further datum belongs here, because it bears on Remark~\ref{I:rem:ghost} and on the
statistics of arithmetic spectra generally. The graph $X^{5,13}$ has $2184$ vertices but
only $55$ distinct eigenvalues, of maximal multiplicity $96$, whereas a random
$3$-regular graph on $2000$ vertices has $2000$ distinct eigenvalues. The degeneracy is
forced by the $\mathrm{PGL}_2(\F_{13})$ representation theory through which the Hecke
operator factors, and it is why the empirical spectral distribution of an LPS graph is a
coarse staircase whose $\mathrm{KS}$ distance decays more slowly than that of a random
graph of half the size. It is also the mechanism behind the Poisson --- not
Gaussian-unitary --- level statistics observed in arithmetic quantum chaos: large Hecke
symmetry produces large degeneracy, and degeneracy is the opposite of level repulsion.
\end{remark}

\subsection{The algorithm, and its verification}

\begin{algorithm}[Hecke $\to$ torsion]\label{I:alg:main}
\textbf{Input:} a prime $p$ (so $q=p$), a torsion-free cocompact $\Gamma\le\PGL_2(\Q_p)$
presented by its quotient graph $Y=\Gamma\backslash X$, i.e. by the integer matrix
$T_p\in M_n(\Z)$.
\begin{enumerate}
\item[(1)] Form $\Delta_p:=(q+1)I_n-T_p\in M_n(\Z)$.
\item[(2)] Compute the Smith normal form over $\Z$,
$U\Delta_pW=\mathrm{diag}(d_1,\dots,d_{n-1},0)$ with $d_1\mid\dots\mid d_{n-1}$
(cost $O(n^3)$ ring operations with the usual bit-length control).
\item[(3)] Output $K_0=\Z\oplus\bigoplus_i\Z/d_i$, $K_1=\Z$, and
$|\Tors K_0|=\prod_id_i=\kappa(Y)$.
\item[(4)] \emph{Certificate.} Independently compute
$\frac1n\prod_{\lambda\neq q+1}(q+1-\lambda)$ from the characteristic polynomial of
$T_p$, and $\frac1n\prod_f(p+1-a_p(f))$ from the Hecke eigenvalues; the three numbers
must agree. Optionally verify Bass's identity \eqref{I:eq:bass} symbolically as a check on
the graph data.
\end{enumerate}
\textbf{Output:} the torsion subgroup of $K_0$ as an explicit finite abelian group,
together with its factorisation into Frobenius traces.
\end{algorithm}

\begin{table}[t]
\centering\small
\renewcommand{\arraystretch}{1.15}\setlength{\tabcolsep}{5pt}
\begin{tabular}{lccccc c}
\toprule
$Y$ & $q$ & $n$ & $\Jac(Y)=\Tors\coker\Delta_p$ & $\kappa(Y)$ & boundary $\Tors K_0$ & $r-1$\\
\midrule
$K_4$              &2& 4&$(\Z/4)^2$                   &$16$    &$\Z/2$&2\\
$K_5$              &3& 5&$(\Z/5)^3$                   &$125$   &$\Z/5$&5\\
$K_6$              &4& 6&$(\Z/6)^4$                   &$1296$  &$\Z/9$&9\\
$K_{3,3}$          &2& 6&$\Z/3\oplus\Z/3\oplus\Z/9$   &$81$    &$\Z/3$&3\\
$Q_3$ (cube)       &2& 8&$\Z/2\oplus\Z/8\oplus\Z/24$  &$384$   &$\Z/4$&4\\
prism $C_4\times K_2$&2& 8&$\Z/2\oplus\Z/8\oplus\Z/24$&$384$   &$\Z/4$&4\\
Wagner $V_8$       &2& 8&$\Z/7\oplus\Z/56$            &$392$   &$\Z/4$&4\\
prism $C_5\times K_2$&2&10&$\Z/19\oplus\Z/95$         &$1805$  &$\Z/5$&5\\
M\"obius $V_{10}$  &2&10&$\Z/11\oplus\Z/165$          &$1815$  &$\Z/5$&5\\
Petersen           &2&10&$\Z/2\oplus(\Z/10)^3$        &$2000$  &$\Z/5$&5\\
Heawood            &2&14&$(\Z/7)^4\oplus\Z/21$        &$50421$ &$\Z/7$&7\\
\bottomrule
\end{tabular}
\medskip
\caption{Exact Smith-normal-form computations. Columns 4--5: the Laplacian channel,
Theorem~\ref{I:thm:structure}; the order always equals
$\frac1n\prod_{\lambda\neq q+1}(q+1-\lambda)$. Columns 6--7: the boundary channel,
Theorem~\ref{I:thm:blind}; the torsion is always cyclic of order $r-1=(q-1)n/2$,
independently of the spectrum, confirming the proof given there. Bass's identity
\eqref{I:eq:bass} was verified symbolically in the first five rows.}
\label{I:tab:snf}
\end{table}

\begin{example}[Level $37$, $p=2$: two elliptic curves, one torsion group]\label{I:ex:37}
The space $S_2(\Gamma_0(37))$ is two-dimensional, spanned by the newforms of the
elliptic curves $37a: y^2+y=x^3-x$ (rank one) and $37b: y^2+y=x^3+x^2-23x-50$ (rank
zero), with $a_2(37a)=-2$ and $a_2(37b)=0$. Correspondingly
\[
\#E_{37a}(\F_2)=2+1-(-2)=5,\qquad \#E_{37b}(\F_2)=2+1-0=3,
\]
and the Hecke--Laplacian of the associated $3$-regular Hecke system has non-zero
eigenvalues $5$ and $3$, so \eqref{I:eq:torsionorder} gives
$|\Tors K_0|=\tfrac13\cdot5\cdot3=5$. The Frobenius traces of two distinct elliptic
curves precipitate into a single cyclic group of order five.
\end{example}

\begin{example}[Level $11$, $p=3$]\label{I:ex:11}
For the definite quaternion algebra of discriminant $11$ the Brandt matrix at $p=3$ is
$B(3)=\begin{psmallmatrix}2&2\\3&1\end{psmallmatrix}$, with spectrum $\{4,-1\}$: the
Eisenstein eigenvalue $q+1=4$ and $a_3(11a)=-1$, matching
$\eta(z)^2\eta(11z)^2=q-2q^2-q^3+2q^4+\cdots$. Then
$q+1-a_3=5=\#E_{11a}(\F_3)=|E_{11a}(\Q)_{\mathrm{tors}}|$.
\end{example}

\begin{remark}[Weighted quotients]\label{I:rem:weights}
Examples~\ref{I:ex:37} and \ref{I:ex:11} are stated at the level of the eigenvalue identity
$q+1-a_p=\#E(\F_p)$, which is unconditional. The passage to the critical group in
Theorem~\ref{I:thm:structure}(iii) requires $\Gamma$ torsion-free, so that $Y$ is a genuine
$(q+1)$-regular graph; Brandt and supersingular isogeny graphs are quotients by groups
with torsion, i.e. graphs of groups with vertex weights $w_i=|\Aut|/2$, and for them
Theorem~\ref{I:thm:structure} holds in its weighted form, with $\Delta_p$ the weighted
Laplacian $w_iB_{ij}=w_jB_{ji}$ and Kirchhoff's theorem replaced by its weighted
version. Congruence subgroups of sufficiently high level are torsion-free and remove the
issue at the cost of enlarging $n$.

The structurally correct home for the weighted case is orbifold, i.e. groupoid,
$K$-theory. A graph of groups $(Y,\Gamma_v,\Gamma_e)$ in the sense of Bass--Serre is the
same datum as an \'etale groupoid $\cG_{Y}$ with unit space $V_Y$ and isotropy
$\Gamma_v$; the object to compute is $K_*(C^*(\cG_Y))$, and the Bass--Serre
Mayer--Vietoris sequence
\[
\cdots\to\bigoplus_{e\in E_Y}K_*\bigl(C^*(\Gamma_e)\bigr)
\xrightarrow{\ \partial\ }\bigoplus_{v\in V_Y}K_*\bigl(C^*(\Gamma_v)\bigr)
\to K_*\bigl(C^*(\cG_Y)\bigr)\to\cdots
\]
replaces $\Z^{E_Y}\to\Z^{V_Y}$ by a sequence of representation rings
$R(\Gamma_v)=K_0(C^*(\Gamma_v))$, of ranks $|\Irr\Gamma_v|$. The vertex weights are then
not an ad hoc normalisation but the ranks of these isotropy contributions: the composite
$\partial\partial^{\,t}$ is the weighted Laplacian $w_iB_{ij}=w_jB_{ji}$ with
$w_i=|\Aut|/2$, the Kirchhoff count acquires the factor $\prod_vw_v^{-1}$ of the
Eichler mass formula, and the apparently fractional indices --- the value
$\tfrac12(4-a_3)$ of Example~\ref{I:ex:11} is not an integer --- become integral once read
in the equivariant group of Proposition~\ref{I:prop:equivariant} rather than in $\Z$. The
$\GL_2$ analogue of Paper XXXVIII's entanglement law belongs here, and we do not develop
it in this paper.
\end{remark}

\subsection{The equivariant refinement}

\begin{proposition}[Equivariant Hecke $K$-theory]\label{I:prop:equivariant}
Let $C$ be a finite group acting on $Y$ by graph automorphisms commuting with $T_p$ ---
in the arithmetic case, $C=\Gal(H/K)$ acting through the class group, as in Paper
XXXVIII. Then $K_0^C(A)=R(C)^{\,n_C}$ for the orbit decomposition, the index map of
Theorem~\ref{I:thm:pimsner} is the $R(C)$-linear extension of $\Delta_p$, and
\[
K_0^C(\cO)\cong\coker_{R(C)}\bigl(\Delta_p\otimes_\Z R(C)\bigr),
\qquad
K_1^C(\cO)\cong\ker_{R(C)}\bigl(\Delta_p\otimes_\Z R(C)\bigr).
\]
The entanglement law of Paper XXXVIII transfers verbatim: free $\Z[C]$-strata of the
vertex module contribute once by Shapiro's lemma, while strata with trivial or sign
action are doubled or sign-twisted by $R(C)$; the orders of the invariant factors are
unchanged, only their representation labels.
\end{proposition}

\subsection{What the two channels say about Langlands}

We summarise the bridge as a table of correspondences, each entry of which is a theorem
or a definition above, not an analogy.

\begin{center}\small
\renewcommand{\arraystretch}{1.2}
\begin{tabular}{lll}
\toprule
Automorphic side & Operator-algebraic side & Reference\\
\midrule
unramified $\pi_p$, Satake $(\alpha_{p,1},\alpha_{p,2})$
& Radon--Nikodym cocycle $e^{-\beta c_\pi}$ & Def.~\ref{I:def:dynamics}\\
local $L$-factor $L_p(\beta,\pi_p)$
& partition function $Z_p(\beta;\pi)$ & Thm~\ref{I:thm:partition}\\
largest pole of $L_p$
& critical inverse temperature $\beta_c(\pi_p)$ & Thm~\ref{I:thm:partition}\\
temperedness (Ramanujan) at $p$
& $\beta_c(\pi_p)=w/2$ & Thm~\ref{I:thm:partition}\\
Eisenstein locus
& the only locus carrying $\KMS$ \emph{states} & Thm~\ref{I:thm:perron}\\
Rankin--Selberg $L(s,\pi\times\widetilde{\pi'})$
& Kakutani/Hellinger obstruction $\Xi_\beta$ & Thm~\ref{I:thm:kakutani}\\
strong multiplicity one
& mutual singularity of phases & Cor.~\ref{I:cor:multone}\\
motivic weight $w$
& width of the equilibrium/obstruction gap & Thm~\ref{I:thm:gap}\\
Hecke recursion $T T_n=T_{n+1}+qT_{n-1}$
& Ihara module $\Z[u^{\pm}]^n/(1-uT_p+qu^2)$ & Def.~\ref{I:def:ihara}\\
Ihara zeta $Z_Y(u)$
& Fitting ideal of the Ihara module & Prop.~\ref{I:prop:bass}\\
Frobenius trace $a_p(f)$, $\#E_f(\F_p)$
& invariant factors of $\SNF(\Delta_p)$ & Thm~\ref{I:thm:structure}\\
Euler characteristic of $Y$ (no arithmetic)
& torsion of the boundary quotient & Thm~\ref{I:thm:blind}\\
\bottomrule
\end{tabular}
\end{center}

\medskip
The reading of the table is the thesis of the paper. Along the top half --- the
thermodynamic rows --- the correspondence is exact but weight-zero: it reproduces class
field theory and Artin data, and Theorem~\ref{I:thm:gap} measures precisely how far it
falls short above weight zero. Along the bottom half --- the $K$-theoretic rows --- the
correspondence is insensitive to temperature and to positivity, and it does carry
weight-two data, but only after the specialisation $u=1$ that separates the Laplacian
channel from the boundary channel. An algorithmic non-commutative class field theory for
$\GL_2$ is therefore not one bridge but two, and Bass's identity \eqref{I:eq:bass} is the
only thing that connects them.

\section*{Acknowledgement of provenance and caveats}

This is an unrefereed preprint and no priority check against the literature has been
carried out. The following are used as known and are not proved here: the Satake
isomorphism and the Hecke recursion; Ihara's theorem and Bass's determinant identity;
Kirchhoff's matrix-tree theorem and the theory of the critical group; the
Eichler--Shimura and Jacquet--Langlands correspondences; Kakutani's dichotomy for product
measures; Jacquet--Shalika non-vanishing for Rankin--Selberg $L$-functions; Renault's
theory of groupoid C*-algebras and the correspondence between $\KMS$ states and conformal
measures; Cuntz--Krieger and Cuntz--Pimsner $K$-theory and Pimsner's exact sequence. The
computations of Table~\ref{I:tab:snf}, and the spectral and torsion-density computations of
Remark~\ref{I:rem:bsnumerics}, are exact or numerically controlled and were carried out
independently of the statements they confirm. Theorem~\ref{I:thm:bstorsion} rests on
Lyons' spanning-tree asymptotics \cite{Lyons2005}, which we quote and do not prove, and
Remark~\ref{I:rem:bsnumerics} on the Kesten--McKay spectral distribution
\cite{Kesten1959,McKay1981}. The one substantive gap left open is the one isolated in
Remark~\ref{I:rem:noalgebra}: the Laplacian channel is a module and an index, not yet a
$C^*$-algebra.
